\documentclass[12pt]{article}
\usepackage[utf8]{inputenc}
\usepackage[T1]{fontenc}
\usepackage{amsmath,amssymb,amsthm}
\usepackage{geometry}
\geometry{margin=1in}
\usepackage{enumitem}
\setlist{nosep}

\linespread{1.1} % Slightly increased line spacing for better readability
% -------------------- Macros --------------------
\newcommand{\N}{\mathbb{N}}
\newcommand{\Z}{\mathbb{Z}}
\newcommand{\Q}{\mathbb{Q}}
\newcommand{\R}{\mathbb{R}}
\newcommand{\C}{\mathbb{C}}
\newcommand{\K}{\mathbb{K}} % used to denote either R or C
\newcommand{\eps}{\varepsilon}
\newcommand{\dd}{\,\mathrm{d}}

\newcommand{\NumberedDefinition}[3]{% #1 number, #2 title, #3 body
\paragraph*{Definition #1 ( #2 ).} #3\par}
\newcommand{\NumberedTheorem}[3]{% #1 number, #2 title, #3 body
\paragraph*{Theorem #1 ( #2 ).} #3\par}
\newcommand{\NumberedProposition}[3]{% #1 number, #2 title, #3 body
\paragraph*{Proposition #1 ( #2 ).} #3\par}
\newcommand{\NumberedLemma}[3]{% #1 number, #2 title, #3 body
\paragraph*{Lemma #1 ( #2 ).} #3\par}
\newcommand{\NumberedCorollary}[3]{% #1 number, #2 title, #3 body
\paragraph*{Corollary #1 ( #2 ).} #3\par}
\newcommand{\NumberedRemark}[3]{% #1 number, #2 title, #3 body
\paragraph*{Remark #1 ( #2 ).} #3\par}
\newcommand{\NumberedExample}[3]{% #1 number, #2 title, #3 body
\paragraph*{Example #1 ( #2 ).} #3\par}


\begin{document}
\section*{3.3. Improper Convergence (Uneigentliche Konvergenz)}

In many cases, it is convenient to extend the set of real numbers by two symbols $+\infty$ (read: plus infinity) and $-\infty$ (read: minus infinity).

\NumberedDefinition{3.3.1}{}{
(1) \textit{We define two symbols $-\infty$ and $+\infty$ and characterize these by the property}
\[
-\infty < x < +\infty \quad \forall x \in \R.
\]
(2) $\sup \emptyset := -\infty$, \textit{with} $\inf \emptyset := +\infty$.

(3) \textit{Let $M \subset \R$ be non-empty. Then we set}
\begin{align*}
\inf M := -\infty \quad &\Leftrightarrow \quad M \textit{ is not bounded below} \\
\sup M := +\infty \quad &\Leftrightarrow \quad M \textit{ is not bounded above.}
\end{align*}
}

\NumberedDefinition{3.3.2}{}{
\textit{We define}
\[
\bar{\R} := \R \cup \{-\infty, \infty\}
\]
\textit{and extend the arithmetic operations on $\bar{\R}$ by}
\begin{align*}
x + \infty &:= +\infty & \forall x &> -\infty \\
x - \infty &:= -\infty & \forall x &< \infty \\
x \cdot (+\infty) &:= \begin{cases} +\infty & x > 0 \\ -\infty & x < 0 \end{cases} \\
x \cdot (-\infty) &:= -x(+\infty) \\
\frac{x}{+\infty} &:= \frac{x}{-\infty} := 0 & \forall x &\in \R \\
\frac{x}{0} &:= \begin{cases} +\infty & x > 0 \\ -\infty & x < 0 \end{cases}
\end{align*}
}

\NumberedRemark{3.3.3}{Suprema/infima in $\overline{\R}$}{Every subset of $\R$ has both an infimum and a supremum \emph{in} $\overline{\R}$. For example, if $A\subset\R$ is unbounded above, then $\sup_{\overline{\R}} A=+\infty$; if $A$ is unbounded below, then $\inf_{\overline{\R}} A=-\infty$.}

\NumberedDefinition{3.3.4}{Limits and cluster points at $\pm\infty$}{Let $(x_n)$ be a sequence in $\R$. We say
\begin{itemize}[leftmargin=*]
	\item $x_n\to +\infty$ (resp. $x_n\to -\infty$) if for every $R\in\R_{\ge0}$ there exists $n_R\in\N$ such that $x_n>R$ (resp. $x_n<-R$) for all $n\ge n_R$.
	\item $a=+\infty$ (resp. $a=-\infty$) is a \emph{cluster point} of $(x_n)$ if for every $R\ge 0$ and every $n\in\N$ there exists $n_R\ge n$ with $x_{n_R}>R$ (resp. $x_{n_R}< -R$).
\end{itemize}}

\NumberedExample{3.3.5}{Improper limits and cluster points}{
\begin{enumerate}[label={(\arabic*)}, leftmargin=*]
	\item $\displaystyle \lim_{n\to\infty} n = +\infty$.
	\item $\displaystyle \lim_{n\to\infty} (-2n) = -\infty$.
	\item The sequence $x_n:=(-1)^n n$ does not converge in $\overline{\R}$, but it has the two cluster points $+\infty$ and $-\infty$ (even indices $\to +\infty$, odd indices $\to -\infty$).
\end{enumerate}}

\NumberedTheorem{3.3.6}{Reciprocals and improper limits}{Let $(x_n)$ be a sequence in $\R$. Then:
\begin{enumerate}[label={(\arabic*)}, leftmargin=*]
	\item If $x_n\to +\infty$ or $x_n\to -\infty$, then $\dfrac{1}{x_n}\to 0$.
	\item If $x_n\to 0$ and $x_n>0$ for all but finitely many $n$, then $\dfrac{1}{x_n}\to +\infty$. If $x_n\to 0$ and $x_n<0$ for all but finitely many $n$, then $\dfrac{1}{x_n}\to -\infty$.
\end{enumerate}}

\paragraph{Proof.}
\begin{enumerate}[label={(\arabic*)}, leftmargin=*]
	\item Suppose $x_n\to +\infty$. Let $\eps>0$. Choose $n_\eps$ with $x_n>1/\eps$ for all $n\ge n_\eps$. Then $0<\tfrac{1}{x_n}<\eps$ for all $n\ge n_\eps$, hence $\tfrac{1}{x_n}\to 0$. The case $x_n\to -\infty$ is analogous using $|1/x_n|$.
	\item Assume $x_n\to 0$ and $x_n>0$ for all but finitely many $n$. Let $R>0$. There exists $n_R$ such that $|x_n|<1/R$ for all $n\ge n_R$. Possibly increasing $n_R$ we may assume $x_n>0$ for all $n\ge n_R$. Then $\tfrac{1}{x_n}>R$ for all $n\ge n_R$, proving $\tfrac{1}{x_n}\to +\infty$. The case with $x_n<0$ eventually is analogous and yields $\tfrac{1}{x_n}\to -\infty$.
\end{enumerate}\qed

\NumberedRemark{3.3.7}{Undefined forms in $\overline{\R}$ and caution}{Using the symbols $+\infty$ and $-\infty$ is often practical, but $\overline{\R}$ is not a field, so several algebraic rules from $\R$ no longer hold. In particular, the following combinations are \emph{not defined}:
\[
	\frac{\infty}{\infty},\ \frac{-\infty}{-\infty},\ \frac{-\infty}{\infty},\ \frac{\infty}{-\infty},\quad
	-\infty + \infty,\ \infty + (-\infty),\quad 0\cdot(+\infty),\ 0\cdot(-\infty).
\]
If such expressions arise during manipulations, one has reached a dead end and should reformulate to avoid them. Moreover, some rules for limits fail in the improper setting. For example, with $a_n=n$ and $b_n=-2n$ we have $a_n\to +\infty$ and $b_n\to -\infty$, but the sum sequence $c_n=a_n+b_n=-n$ converges to $-\infty$, while the `sum of limits' $+\infty+(-\infty)$ is not defined. Similar issues occur for quotients of sequences that tend to $\pm\infty$.}

\NumberedTheorem{3.3.8}{Monotone sequences converge in $\overline{\R}$}{Every monotone increasing (resp. monotone decreasing) sequence $(x_n)_{n\in\N}$ in $\R$ converges in $\overline{\R}$, and
\[
	\lim_{n\to\infty} x_n = \sup\{x_n: n\in\N\}\quad \big(\text{resp. }\ \lim_{n\to\infty} x_n = \inf\{x_n: n\in\N\}\big).
\]
\paragraph{Proof.} Consider the increasing case; the decreasing case is analogous. If $x_n=-\infty$ for all $n$, then clearly $\lim x_n=-\infty$ and $\sup\{x_n:n\in\N\}=-\infty$. Otherwise, choose $n_0\in\N$ with $x_{n_0}>-\infty$ and examine the tail $(x_n)_{n\ge n_0}$. If it is bounded above in $\R$, then by Theorem~3.2.2 it converges in $\R$ to its supremum. If it is not bounded above, then for every $R>0$ there exists $n_R\ge n_0$ with $x_{n_R}>R$, and monotonicity gives $x_n>R$ for all $n\ge n_R$, hence $x_n\to +\infty$. In either case the limit in $\overline{\R}$ equals the supremum. \qed}

\paragraph{Tail suprema and infima.} Let $(x_n)_{n\in\N}$ be a sequence in $\R$ and define, for $n\in\N$,
\[
	y_n := \sup\{ x_k : k\ge n\}, \qquad z_n := \inf\{ x_k : k\ge n\}.
\]
Then $(y_n)$ is monotone decreasing in $\overline{\R}$ and $(z_n)$ is monotone increasing in $\overline{\R}$. By Theorem~3.3.8, both $(y_n)$ and $(z_n)$ converge in $\overline{\R}$.

\NumberedDefinition{3.3.9}{Limsup and liminf}{Let $(x_n)_{n\in\N}$ be a sequence in $\R$. Then
\[
	\limsup_{n\to\infty} x_n := \lim_{n\to\infty} \sup\{x_k : k\ge n\},
	\qquad
	\liminf_{n\to\infty} x_n := \lim_{n\to\infty} \inf\{x_k : k\ge n\}.
\]}

\NumberedRemark{3.3.10}{}{Because $\sup\{x_k : k\ge n\} \ge \inf\{x_k : k\ge n\}$ holds for every $n$, we always have
\[
	\limsup_{n\to\infty} x_n \;\ge\; \liminf_{n\to\infty} x_n.
\]}

\NumberedTheorem{3.3.11}{Limsup/liminf and cluster points}{Let $(x_n)_{n\in\N}$ be a sequence in $\R$. Then $\limsup_{n\to\infty} x_n$ (resp. $\liminf_{n\to\infty} x_n$) is the greatest (resp. smallest) cluster point in $\overline{\R}$ of $(x_n)$.}

\paragraph{Proof.} Set $x := \limsup_{n\to\infty} x_n$ and define $y_n := \sup\{x_k : k\ge n\}$. Then $(y_n)$ is monotone decreasing and $y_n\to x$.

\emph{Case 1: $x=+\infty$.} Since $y_n\to +\infty$ and $(y_n)$ is decreasing, we have $\forall R>0\ \forall n\in\N\ \exists k>n:\ x_k>R$ (because $+\infty$ is the least upper bound of the tail sets $\{x_k:k\ge n\}$). Thus $+\infty$ is a cluster point of $(x_n)$ in $\overline{\R}$, and evidently there is no larger cluster point.

\emph{Case 2: $x\in\R$.} First we show that no cluster point $z$ of $(x_n)$ is greater than $x$. Let $\eps>0$. Since $y_n\downarrow x$, there exists $n_\eps$ with $x\le y_n < x+\eps$ for all $n\ge n_\eps$. Hence
\[
	\#\{k\in\N : x_k \ge x+\eps\} \le n_\eps < \infty.
\]
Therefore every cluster point $z$ satisfies $z\le x+\eps$, and since $\eps>0$ is arbitrary, $z\le x$.

Next we show that $x$ is a cluster point of $(x_n)$. Let $\eps>0$ and $n\in\N$. Since $x\le y_n < x+\eps$ and $y_n$ is the least upper bound of $\{x_k:k\ge n\}$, there exists $k\ge n$ with
\[
	x \le x_k < x+\eps.
\]
Thus $x$ is indeed a cluster point. Hence $x$ is the greatest cluster point.

\emph{Case 3: $x=-\infty$.} Let $R\in\R_{\ge0}$. Then there exists $n_R\in\N$ such that $y_n<-R$ for all $n\ge n_R$. It follows that
\[
	\#\{k\in\N : x_k \ge -R\} \le n_R < \infty.
\]
Therefore any cluster point $z$ must satisfy $z\le -\infty$, i.e. $z=-\infty$. Moreover, for every $R\in\R_{\ge0}$ and every $n\in\N$ there exists $k\ge n$ with $x_k<-R$, proving that $-\infty$ is a cluster point. This shows that the smallest cluster point equals $-\infty$ in this case.

The proof for $\liminf_{n\to\infty} x_n$ is analogous. \qed

\NumberedExample{3.3.12}{}{Let $x_n := (-1)^n\,\dfrac{n}{n+1}$. Then
\[
	\limsup_{n\to\infty} x_n = 1
	\qquad\text{and}\qquad
	\liminf_{n\to\infty} x_n = -1.
\]
Indeed, with $y_n := \sup\{x_k:k\ge n\}$ and $z_n := \inf\{x_k:k\ge n\}$ we have $y_n=1$ for all $n\in\N$ and $z_n=-1$ for all $n\in\N$.}


\end{document}